%\chapter{ABSTRACT}
\clearpage
\phantomsection
\addcontentsline{toc}{chapter}{ABSTRACT}
\begin{center}
%    \textbf{\large{\judulen}}\\[0.5cm]
%    by:\\
%    \penulis\\
%    \nim\\[2em]
    \textbf{ABSTRACT}\\[0.5cm]
\end{center}

The website has evolved into a complex platform, one of which is its use as a dashboard for displaying information from IoT devices. In this context, a dashboard website serves as a visual interface to present data collected by IoT sensors and to visualize that data. To create IoT data visualizations, there are various libraries available, each with its own advantages and disadvantages. Considering factors such as rendering performance, flexibility, and its implementation, JavaScript libraries are a suitable choice because of their flexibility and compatibility with all modern browsers without requiring additional dependencies. This study compares three JavaScript libraries, Chart.js, D3.js, and Highcharts, to help developers choose the most appropriate one for handling continuously updated IoT data. The evaluation of these JavaScript visualization libraries focuses on real-time data rendering performance, chart scalability, and the efficiency of CPU and memory usage. The results show that Chart.js delivers good rendering performance, is efficient in memory and CPU usage, and maintains high stability without significant fluctuations. D3.js excels in customization but is less stable for large datasets and tends to consume more memory and CPU resources. Highcharts maintains rendering speed, stability, and memory efficiency under heavy data loads because of its built-in optimization features, making it suitable for large-scale IoT visualizations. It is hoped that the findings of this study will provide insights and guidance in selecting an optimal visualization library to improve dashboard performance and overall system efficiency.


\noindent Key words: \keywords